% Auto-generated from references/originals/welcome-to-ucas-original.pdf by scripts/bootstrap_sources.py.
% The extraction is intentionally conservative; revise wording and structure during editing.
\chapter{来到国科大的第一课}
\sourceattribution{本章逆向整理自原 PDF《来到国科大的第一课！》，原作者：试管，中国科学院大学本科部 23 级学生。}

来到国科大的第一课！欢迎新同学就读中国科学院大学，加入我们安贫乐道的队伍！面对崭新的大学生活，每位同学都会感到陌生、期待与一丝担忧。

请放心，新同学无论提出什么问题，都会有热情的学长学姐尽可能地帮助你解答，协助你度过“新手保护”阶段，顺利适应玉泉路的校园生活。

每一所大学的学习生活细节都不是三言两语能够讲完的，不过你可以通过随时询问来获得帮助。但在此之前，了解以下国科大的特色文化，有助于你更好融入国科大的生活，也可以让为你解答问题的老东西们轻松不少！

\section{别急}

来到国科大的第一课是“别着急”。我校行政效率较低，部门之间权责归属不够明确，本科部行政老师身兼多职分身乏术，这是国科大的特色。任何与我校相关的事情，要做好最后一刻才接到通知的准备，要做好啥事都别急的心理预期，要了解到即使是通知已经下达了有可能还有老师/同学不知道的悲惨事实。坚持“我不着急，因为急也没用，我不急有的是人急，别人不急我为什么要急”的思想原则，可以让你的大学生活轻松不少。

（期待我继续写？您也别急）

\section{有事先上网}

某些问题是在门户网站（如百度、必应）以及

我校常见网站（如学校官网、本科教育网、迎新服务网、教务处网站）就能解决的，那么不需要麻烦你的学长学姐帮你回答（可以提前熟悉一下我校常见网站的功能）

；某些问题你可能有印象“我之前在群里好像看到过”，那么麻烦同学使用聊天记录搜索而不是反复提问。等你度过了“新手保护期”，可能有前辈会在交流技术问题的时候告诉你 RTFM（去读手册）或者 STFW（去上网查）

，其实是一样的意思。

\section{“科气”是坏文明}

刚刚结束高考的同学们，在应试教育中可

能会出现“羡慕……”“膜拜……”“我好菜”“我高中的时候学过……”之类的心理或者言论。在国科大，这是被广泛排斥的。国科大的生活不是只有什么做题、什么绩点，我们也在尽全力避免这所学校的内卷化、高中化，这是我们和某兄弟院校最大的区别。在新手保护期，我们对这些言论只会提出温柔的警告。但是如有再犯，或者开学后两三个月仍出现这类症状，我建议移送中科大，并请你不要介意我们用各种词汇来攻击发言者（当然，我们也会有群专门收容 GPA 爱好者）

。我也劝诸位同学在这个暑假别搞什么“预习”“抢跑”，原因有二：一则这个暑假可能是你享受到的最后一个暑假，此后就要忍受假期也要被安排事情；二则所谓的“预习”可能完全偏离课程，或者预习两个月的结果老师一节课就讲完了。此外，请不要在公开场合收集课程考试的往年题目。

\section{拥有最基本的社会素质与道德修养}

如果有同学解答了你的

问题，或者请你提供更多信息细化你的问题，请不要默不作声、啥话不说；如果在国科大社群有任何不愉快的线上经历，及时解释/道歉即可，毕竟你目前拥有新手保护。

以我个人目前的经历来看，只要牢记以上四点，在新手保护期内，你是百无禁忌、可以享受到学长学姐的无私帮助的。请尽情享受你的这个暑假，期待你能享受在国科大安贫乐道四年的日子！
